% Created 2026-03-18 Wed 05:54
% Intended LaTeX compiler: pdflatex
\documentclass[11pt]{article}
\usepackage[utf8]{inputenc}
\usepackage[T1]{fontenc}
\usepackage{graphicx}
\usepackage{longtable}
\usepackage{wrapfig}
\usepackage{rotating}
\usepackage[normalem]{ulem}
\usepackage{amsmath}
\usepackage{amssymb}
\usepackage{capt-of}
\usepackage{hyperref}
\author{Chris Hurst}
\date{2026-03-13}
\title{Principal Component Analysis of a Constant Position Kalman Filter}
\hypersetup{
 pdfauthor={Chris Hurst},
 pdftitle={Principal Component Analysis of a Constant Position Kalman Filter},
 pdfkeywords={},
 pdfsubject={},
 pdfcreator={Emacs 29.3 (Org mode 9.6.15)}, 
 pdflang={English}}
\begin{document}

\maketitle

\section*{Introduction}
\label{sec:org621ab24}
This is a short investigation on the principal components of
\begin{enumerate}
\item Kalman filter posterior states
\item Kalman filter residuals
\end{enumerate}
for a constant position Kalman filter.

\section*{Analysis}
\label{sec:orgd05b8df}
First, let's look at an object not drifting which has random noise about (0,0).

\begin{figure}[htbp]
\centering
\includegraphics[width=.9\linewidth]{/home/chris/org/test-blog/org/static/images/random_noise.png}
\caption{\label{fig:orgce941f8}Random noise object.  This is how it \textbf{should} behave.  The Kalman filter slightly lags behind the measurements.}
\end{figure}

The bottom subplot shows the residuals and the direction of the principal eigenvector.
The eigenvalue ratio of the residuals is 3.45.  


Now an object with drift, and little noise.

\begin{figure}[htbp]
\centering
\includegraphics[width=.9\linewidth]{/home/chris/org/test-blog/org/static/images/drift_smooth.png}
\caption{\label{fig:orgeba0176}Drifting object with little noise.  The principal eigenvector for both residuals shows we get the direction of drift.}
\end{figure}

Here the eigenvectors are pointing in the direction of drift.
There is little variance in the perpendicular direction of the object's motion.
The eigenvalue ratio of the residuals is 582.7. Quite large.

But what about the case where an object is drifting but has fluctuations perpendicular to the motion of travel?

\begin{figure}[htbp]
\centering
\includegraphics[width=.9\linewidth]{/home/chris/org/test-blog/org/static/images/drift_with_up_and_down.png}
\caption{\label{fig:org166e22d}Drifting object with high variance in the perpendicular direction of travel.}
\end{figure}

The principal eigenvector for the residuals points in the direction of most variance, \emph{not} the direction of drift.
This is not quite what I expected before.
The eigenvalue ratio is 10.5.


Additionally, this means that if an object has high variance but zero mean, we also get misleading results.

\begin{figure}[htbp]
\centering
\includegraphics[width=.9\linewidth]{/home/chris/org/test-blog/org/static/images/high_variance_y.png}
\caption{\label{fig:orgc835cb1}Object with no drift, but measurements bouncing up and down between y, with little change in x.}
\end{figure}

Again, misleading.
The eigenvector ratio is 34445.5, which is way too much.

\section*{A new idea}
\label{sec:org47bbee5}
So it looks like some fancy EVD isn't going to work for these purposes.
Let's try something else.

The new plan is to take the mean of the residuals and then normalize them against their covariance.
\begin{equation}
\tilde{S} = \frac{1}{N-1} \nu^{\text{T}} \nu
\end{equation}
Perform EVD on the \(2x2\) matrix \(\tilde{S}\) to find \(\tilde{S}^{-1/2}\).
I am using \(\tilde{S}\) instead of \(S\) to distinguish between the innovation covariance of the kalman filter.
This is a `sample' innovation covariance.
\begin{equation}
\tilde{S}^{-1/2} = \Lambda D^{-1/2} \Lambda^{\text{T}}
\end{equation}
And now the normalized innovations are
\begin{equation}
\tilde{\nu}_i = \tilde{S}^{-1/2} \nu_i
\end{equation}
And the drift vector can be found by taking the average of \textasciitilde{}\{\(\nu\)\}\_i:
\begin{equation}
\bar{d} = \frac{1}{N} \sum_i^N{\tilde{\nu}_i}
\end{equation}

Lets plot this using the same scenarios as before.
The bottom plots are the average drift vector.

\begin{figure}[htbp]
\centering
\includegraphics[width=.9\linewidth]{/home/chris/org/test-blog/org/static/images/drift_random_noise.png}
\caption{\label{fig:orge423b9c}Random noise object.  The drift magnitude is \(d=0.109\).}
\end{figure}


\begin{figure}[htbp]
\centering
\includegraphics[width=.9\linewidth]{/home/chris/org/test-blog/org/static/images/drift_drift_smooth.png}
\caption{\label{fig:org63905b7}Drifting object with little noise. The drift magnitude is \(d=0.912\).}
\end{figure}


\begin{figure}[htbp]
\centering
\includegraphics[width=.9\linewidth]{/home/chris/org/test-blog/org/static/images/drift_drift_with_up_and_down.png}
\caption{\label{fig:orge2e329c}Drifting object with high variance in the perpendicular direction of travel. The drift magnitude is \(d=0.907\).}
\end{figure}


\begin{figure}[htbp]
\centering
\includegraphics[width=.9\linewidth]{/home/chris/org/test-blog/org/static/images/drift_high_variance_y.png}
\caption{\label{fig:org42d344a}Object with no drift, but measurements bouncing up and down between y, with little change in x. The drift magnitude is \(d=0.205\).}
\end{figure}

This is a lot more in line with a metric for something that is drifting.
\end{document}
